\section{Related Work}
\label{sec:related}

\subsection{Vulnerability Reproducibility and PoC Generation}

Mu et al. study the reproducibility of crowd-reported vulnerabilities
through structured analyst experiments and a separate practitioner
survey~\cite{mu2018crowdrepro}. Their analysis connects reproduction
outcomes to available information, information recovery, and debugging.
It establishes a useful precedent for studying the workflow around a
vulnerability report, including the distinction between frequently
missing information and information consequential to reproduction. Our
setting replaces the analyst workflow with general LLM agents and focuses
on binary firmware, execution evidence, and subsequent agent behavior.

Fang et al. investigate exploitation of one-day vulnerabilities with
different information conditions~\cite{fang2024oneday}. Zhao et al. study
Web PoC generation, decompose failures, and compare code context and
feedback strategies~\cite{zhao2025webpocstudy}. Patch-to-PoC examines
agentic reproduction of Linux kernel vulnerabilities using patch input
and dedicated execution support~\cite{krepro}. These studies show how
provided information and tool conditions affect the task being evaluated.
Their prepared environments and success definitions differ from the
CVE-initiated firmware workflow here, so their rates are not direct
performance baselines for our agents.

CVE-GENIE automatically gathers CVE resources, reconstructs vulnerable
software environments, and develops verifiable exploits~\cite{cvegenie}.
SEC-bench likewise uses automated construction and validation for software
security tasks~\cite{secbench}. These systems explicitly address
environment preparation; we do not claim that prior research uniformly
hands agents a ready target. Our contribution is intended to be a
process-level empirical account of general coding agents on selected IoT
firmware cases, including the relationship between target preparation,
target changes, and outcome reporting.

Firmware rehosting research studies the conditions required to execute
embedded software away from its device. Firmadyne and FirmAE develop
automated approaches; Greenhouse focuses on single-service rehosting;
Pandawan examines progress in Linux-based firmware
rehosting~\cite{firmadyne,firmae,greenhouse,pandawan}. The challenges
survey systematizes hardware and software dependencies~\cite{challenges}.
FirmAgent combines fuzzing and LLM analysis for IoT vulnerability
discovery~\cite{firmagent}. These works provide technical context for our
obstacle analysis. Evaluating general agents' handling of those conditions
is distinct from comparing the underlying rehosting tools themselves.

\subsection{Empirical Studies of Security Agents}

Happe and Cito combine quantitative evaluation and expert analysis of
agent trajectories in an enterprise network~\cite{happe2026enterprisenetworks}.
Shao et al. investigate offensive-security challenge solving across
interactive and automated workflows~\cite{shao2024offensivestudy}.
Their evidence motivates studying commands, feedback, and recovery
alongside completion, while their environments and degrees of human
participation differ from our task.

Huang et al. relate functional properties of agent architectures to
penetration-testing performance, failure modes, and architectural
interventions~\cite{huang2025capabilitiesperformance}. Probst et al.
derive candidate enhancements from observed privilege-escalation failures
and test their combinations~\cite{probst2026localprivesc}. Incalmo uses
failure analysis to motivate a new multi-host red-teaming
system~\cite{singer2026incalmo}. These papers illustrate an observation-to-
intervention or observation-to-system structure. Our present analysis
does not claim equivalent intervention evidence or a new agent architecture.

Yu et al. decouple reconnaissance and exploitation by supplying missing
context~\cite{yu2026decoupling}. Erdem examines repeated attack runs on
a fixed target and discusses service errors and other execution
confounds~\cite{erdem2026consistency}. Garc\'{i}a-Carrasco et al. study
local model configurations on a small CyberGym subset using repeated runs
and trajectory examples~\cite{garcia2026local}. Together, these studies
inform our treatment of prerequisite failures, repeated trials, and
configuration-dependent performance. Each provides evidence within its
own case and tool conditions rather than a universal ordering of models.

The compared studies assign different target-preparation responsibilities.
A prepared runtime can still expose a difficult exploitation problem,
while a CVE-only task can be easier when a reusable public recipe exists.
Initial information and execution responsibilities must be considered
together; their success rates are not interchangeable.


The comparison also identifies the specific gap our process analysis is
intended to address. Human reproduction studies explain difficulties of
recovering missing information. Prepared-target studies can investigate
payload construction and interaction after initialization. Specialized
construction systems can demonstrate that their pipeline prepares
particular software environments. Our observations concern how general
agents navigate the combined firmware workflow and what happens when
their tested object or entry path changes during that attempt. These are
complementary questions, and a difference in starting condition alone is
not evidence of a stronger or weaker agent.

Measurement validity is also central in adjacent tasks. SecLLMHolmes
studies stability and reasoning-related outputs in vulnerability
identification~\cite{ullah2024secllmholmes}; Risse and B\"ohme examine
whether vulnerability detectors remain reliable under targeted
evaluation changes~\cite{risse2024limits}. CodeRover-S distinguishes
plausible patches from manually assessed correct patches and practical
integration~\cite{zhang2026ossfuzzrepair}. Our distinction between a
working demonstration and evidence on a specified original target shares
this concern with what a reported outcome actually measures.

\subsection{Evaluation Resources for Security Agents}

CVE-Bench supplies real-world Web targets and explicit attack
objectives~\cite{cvebench}. CyberGym provides prepared reproduction
environments for real software vulnerabilities~\cite{cybergym}.
BountyBench organizes detection, exploitation, and patching around
real-world systems and vulnerability instances~\cite{bountybench}.
These resources enable controlled comparisons, but differ in initial
information, environment preparation, intended outcome, and verification.
The survey of offensive-agent evaluation practices highlights such design
choices and the difficulty of reducing long workflows to linear
progress~\cite{happe2025offensivepractices}.

\sys{} serves as the resource behind the present study. Its relevance
rests on traceable observations of the selected workflow. We make neither
a largest-dataset claim nor a claim that external benchmarks cannot
validate original targets. The empirical comparison concerns what is
observed when agents themselves confront the particular target-preparation
requirements and evidence boundaries described in this paper.
